\documentclass[11pt]{article}
\usepackage{style}
\geometry{
    a4paper,
    left=1in,
    right=1in,
    top=1in,
    bottom=1in
}
\setstretch{1.1}
\setlength{\parskip}{1.2em}
\setlength{\emergencystretch}{2em}

% Sans-serif Style
\setmainfont[
    Path=fonts/tex-gyre/,
    Extension=.otf,
    UprightFont=texgyreheros-regular,
    BoldFont=texgyreheros-bold,
    ItalicFont=texgyreheros-italic,
    BoldItalicFont=texgyreheros-bolditalic
]{texgyreheros}

\setCJKmainfont[
    Path=fonts/cjk/,
    Extension=.ttf,
    UprightFont=NotoSansSC-Regular,
    BoldFont=NotoSansSC-SemiBold,
    ItalicFont=ChillKai,
    BoldItalicFont=ChillKai
]{Mix-Sans}

% Serif Style
% \setmainfont[
%     Path=fonts/tex-gyre/,
%     Extension=.otf,
%     UprightFont=texgyretermes-regular,
%     BoldFont=texgyretermes-bold,
%     ItalicFont=texgyretermes-italic,
%     BoldItalicFont=texgyretermes-bolditalic
% ]{texgyretermes}

% \setCJKmainfont[
%     Path=fonts/cjk/,
%     Extension=.ttf,
%     UprightFont=NotoSerifSC-Regular,
%     BoldFont=NotoSerifSC-SemiBold,
%     ItalicFont=ChillKai,
%     BoldItalicFont=ChillKai
% ]{Mix-Serif}

\setCJKmonofont[
    Path=fonts/cjk/,
    Extension=.ttf,
    UprightFont=NotoSansSC-Regular,
]{NotoSansSC}

\setmathfont{STIX Two Math}

\setmonofont[
    Path=fonts/JetBrainsMono/,
    Extension=.ttf,
    UprightFont=JetBrainsMono-Regular,
    BoldFont=JetBrainsMono-Bold,
    ItalicFont=JetBrainsMono-Italic,
    BoldItalicFont=JetBrainsMono-BoldItalic
]{JetBrainsMono}

% 图片占位：figs/ 下对应文件存在时自动渲染真实图片，否则显示灰色占位框
\newcommand{\placeholderfig}[2]{%
    \IfFileExists{#1}%
        {\includegraphics[width=#2\textwidth]{#1}}%
        {\textcolor{gray!40}{\rule{#2\textwidth}{6cm}}}%
}

\frontpagestyle{t1}

\reportlogo{figs/HUST-logo.pdf}
\reporttitle{Awesome-Micro-Diffusion-Model：\\从零开始的迷你扩散模型}
\reportauthor{张禹阳}
\reportaffiliation{电子信息与通信学院}
\reportcontact{yuyangzh@hust.edu.cn}
\reportabstract{本报告记录在 MNIST 数据集上从零实现并训练迷你扩散模型（DDPM）的完整过程。去噪网络为参数量受限（$\leq$5M，实际 1.85M）的 U-Net，全部代码仅使用原生 PyTorch 实现，在纯 CPU 的约束下完成训练，从标准正态噪声经 ancestral sampling 生成 8$\times$8 样本网格。报告内容：前向闭式加噪与反向去噪后验的推导、U-Net 与训练设计、EMA 权重指数滑动平均；两项扩展——基于 DDIM 的不同步数采样（$\eta$ 参数化，$\eta=0$ 时确定性采样）与基于 classifier-free guidance 的指定数字生成；以及训练耗时超标（约 6h 优化至 1h 量级）、$\eta=0$ 确定性采样发散（定位为余弦调度末端 $\hat{x}_0$ 反解的数值放大，以截断修复）等工程问题的分析与解决。}

% ------------------------------------------- %
\addbibresource{refs.bib}
% ------------------------------------------- %

\begin{document}

\makereporttitle
\newpage
\tableofcontents
\newpage

\section{实验概述}

一位画家，从一幅完工的画开始观察：\textbf{前向}过程向这幅画泼五颜六色的颜料，一步步破坏它；\textbf{反向}过程从一片狼藉中，一步步去掉泼上去的颜料。扩散模型正是这一直觉的数学化：前向逐步向图像注入高斯噪声，直到变成纯噪声；反向训练一个去噪网络，从纯噪声中一步步把图像还原出来。

本实验在 MNIST 数据集~\parencite{lecun1998gradient}（训练集 60,000 张、测试集 10,000 张，每张 $28\times28\times1$）上从零实现并训练一个 DDPM~\parencite{ho2020denoising}，全部代码仅使用原生 PyTorch~\parencite{paszke2019pytorch}，不借助 diffusers 等高级库。实验要求与完成情况：

\begin{itemize}
    \item \textbf{前向加噪}：噪声调度（线性/余弦两种），一步到位的闭式加噪；
    \item \textbf{去噪网络}：U-Net（参数量 $\leq 5$M，实际 1.85M），以 MSE 预测噪声并训练；
    \item \textbf{反向采样}：ancestral sampling，从纯噪声逐步去噪，输出 8$\times$8 样本网格；
    \item \textbf{约束}：CPU 训练在 2h 水平（本机为 AMD Ryzen7 5800H），支持 CPU/GPU 自动检测；
    \item \textbf{扩展一}：不同采样步数对生成质量与耗时的影响（DDIM 实现）；
    \item \textbf{扩展二}：指定数字采样（条件扩散 + classifier-free guidance）。
\end{itemize}

\section{数学原理}

本节按"前向加噪 $\rightarrow$ 反向去噪 $\rightarrow$ 去噪网络 $\rightarrow$ 训练技巧（EMA）$\rightarrow$ 两项扩展"的顺序给出核心推导。完整推导见项目 \texttt{README.md}，此处为简要分析。

\subsection{前向加噪：从逐步加噪到一步到位}

原始图像 $x_0$ 是高维空间中的一个向量（MNIST 即 784 维）。总步数 $T$，每一步由上一步的图像生成本步的带噪图像：

\begin{equation}
    \boxed{x_{t} = x_{t-1}\sqrt{1-\beta_t} + \epsilon_t \sqrt{\beta_t}}
\end{equation}

噪声向量 $\epsilon_t$ 的每个元素独立地从 $\mathcal{N}(0,1)$ 抽样；$\beta_t$ 是每步的噪声方差，按预先设计的调度变化。本实现支持两种调度：

\begin{itemize}
    \item \textbf{线性调度}~\parencite{ho2020denoising}：$\beta_1=10^{-4},\ \beta_T=0.02$，中间线性插值；
    \item \textbf{余弦调度}~\parencite{nichol2021improved}：$\bar{\alpha}_t = f(t)/f(0)$，$f(t)=\cos^2\!\big(\frac{t/T+s}{1+s}\cdot\frac{\pi}{2}\big)$，$s=0.008$，$\beta_t = 1-\bar{\alpha}_t/\bar{\alpha}_{t-1}$（截断至 $\leq 0.999$）。
\end{itemize}

逐步加噪可以合并。记 $\alpha_t = 1-\beta_t$，$\bar{\alpha}_t = \prod_{i=1}^t \alpha_i$，利用"独立正态分布之和仍为正态、方差相加"的性质递推，得到一步到位的闭式公式：

\begin{equation}
    \boxed{x_t = x_0\sqrt{\bar{\alpha}_t} + \epsilon \sqrt{1-\bar{\alpha}_t}}
\end{equation}

训练时无需逐步模拟，对每个样本随机采一个 $t$ 即可直接构造 $x_t$。

\subsection{反向去噪：后验分布与 ancestral sampling}

目标是根据 $x_t$ 和 $t$ 求 $x_{t-1}$ 的分布。在给定 $x_0$ 的条件下，$(x_t, x_{t-1})$ 服从二维正态分布；计算协方差并利用"二维正态的条件分布仍是正态"，可得后验 $q(x_{t-1}|x_t,x_0) = \mathcal{N}(\tilde{\mu}_t, \tilde{\beta}_t)$，其中：

\begin{equation}
    \boxed{\tilde{\mu}_t = \frac{1}{\sqrt{\alpha_t}} \left(x_t - \frac{\beta_t}{\sqrt{1-\bar{\alpha}_t}}\epsilon_{\theta}(x_t;t)\right)}
    \qquad
    \boxed{\tilde{\beta}_t = \frac{1-\bar{\alpha}_{t-1}}{1-\bar{\alpha}_t}\,\beta_t}
\end{equation}

真实的 $\epsilon$ 未知，用网络预测 $\epsilon_\theta$ 代替。去噪公式即为 ancestral sampling 的一步：

\begin{equation}
    \boxed{x_{t-1} = \frac{1}{\sqrt{\alpha_t}} \left(x_t - \frac{\beta_t}{\sqrt{1-\bar{\alpha}_t}}\epsilon_{\theta}(x_t;t)\right) + z\sqrt{\tilde{\beta}_t}}
\end{equation}

其中第二项 $z\sqrt{\tilde{\beta}_t}$（$z\sim\mathcal{N}(0,1)$）提供随机性，$t=1$ 时不再加噪声。

\subsection{去噪网络 U-Net}

U-Net~\parencite{ronneberger2015unet} 由编码器、瓶颈、解码器组成，跳跃连接按 channel 维拼接。本项目经 CPU 实测调优后的配置（调优过程见第 4 节）：

\begin{table}[h]
    \centering
    \caption{U-Net 结构（输入 $[B,1,28,28]$，输出同形状预测噪声，参数量 1.85M）}
    \label{tab:unet}
    \begin{tabular}{l l l l}
        \toprule
        \textbf{位置} & \textbf{分辨率} & \textbf{通道数} & \textbf{组成} \\
        \midrule
        输入卷积 & $28^2$ & 24 & $3\times3$ conv \\
        编码器 L0/L1/L2 & $28^2/14^2/7^2$ & 24/48/96 & ResBlock $\times$ 2，MaxPool 降采样 \\
        瓶颈 & $7^2$ & 96 & ResBlock $\rightarrow$ 4 头自注意力 $\rightarrow$ ResBlock \\
        解码器 L2/L1/L0 & $7^2/14^2/28^2$ & 96/48/24 & 最近邻上采样，concat 跳跃连接，ResBlock $\times$ 2 \\
        输出层 & $28^2$ & 1 & $1\times1$ conv 恢复通道数 \\
        \bottomrule
    \end{tabular}
\end{table}

残差块的结构为 $\mathrm{Conv1} \rightarrow \mathrm{GroupNorm} \rightarrow \mathrm{SiLU} \rightarrow \mathrm{Conv2} \rightarrow \mathrm{GroupNorm} \rightarrow \mathrm{SiLU}$，并把输入经跳跃连接加回输出。时间步 $t$ 经正弦嵌入（$f_i = 10000^{-2i/d}$，偶数位 $\sin(t f_i)$、奇数位 $\cos(t f_i)$）加两层 MLP 得到统一的时间向量；每个残差块再拥有专属的两层 MLP，把它投影为本块通道数的 $\gamma$（scale）与 $\beta$（shift），作用于第一个 GroupNorm 之后——时间信息由此注入到每一个残差块。瓶颈处把 $7\times7=49$ 个像素视为 token，按 channel 切 4 头做自注意力。

\subsection{EMA 权重指数滑动平均}

训练中参数 $\theta$ 的每步更新都带有 batch 抽样的随机性，最后时刻的参数只是抖动轨迹上的一个随机快照。EMA 维护一份影子参数，每步做一次插值：

\begin{equation}
    \boxed{\theta_{ema} \leftarrow \rho\,\theta_{ema} + (1-\rho)\,\theta}
\end{equation}

把递推展开，$\theta_{ema}$ 是参数历史轨迹的指数加权平均，有效窗口约 $\frac{1}{1-\rho}$ 步：$\rho=0.999$ 对应约 1000 步（约 2 个 epoch），窗口内的高频抖动被平均掉，相当于对参数轨迹做低通滤波。影子参数不参与梯度；采样时加载 EMA 权重（DDPM 原文亦如此~\parencite{ho2020denoising}）。

\subsection{扩展一：少步采样与 DDIM}
\label{sec:ddim}

ancestral sampling 必须严格走满全部 $T$ 步，推理耗时与 $T$ 成正比；能否不重训、只用 $S \ll T$ 步完成采样？

\textbf{1. 突破口：训练只约束"边缘分布"}。扩散模型的训练损失只约束单点边缘分布 $q(x_t|x_0)=\mathcal{N}(x_0\sqrt{\bar{\alpha}_t},\,1-\bar{\alpha}_t)$——只检查"第 $t$ 步的图长什么样"，至于这张图是从 $t-1$ 步爬楼梯过来的，还是从起点空降过来的，损失函数根本不关心。只要保证边缘分布不变，从 $x_T$ 到 $x_0$ 的"行走路径"就可以重新规划，从而绕开必须走满 $T$ 步的限制。

\textbf{2. 重新设计路径：非马尔可夫前向过程}。DDPM 的前向是马尔可夫链：$x_t$ 只依赖 $x_{t-1}$，必须一阶一阶爬楼梯。DDIM~\parencite{song2021ddim} 构造一族\textbf{非马尔可夫}前向过程，允许 $x_t$ 同时依赖 $x_0$ 与上一步，从而可以直接"跳步"——从 $x_{\tau_i}$ 转移到子序列上的 $x_{\tau_{i-1}}$。直观理解：DDPM 规定必须爬楼梯；DDIM 说只要到达时窗外的风景（边缘分布）与爬楼梯看到的完全一样，坐电梯直达任意楼层就是合法的。

\textbf{3. 反向更新的工程构造}。先用边缘分布反解出预测的干净图像 $\hat{x}_0$；既然 $x_{\tau_{i-1}}$ 必须落在以 $\hat{x}_0$ 为中心、方差为 $1-\bar{\alpha}_{\tau_{i-1}}$ 的边缘分布上，它的结构就必然长成"原图方向 + 噪声方向"：

\begin{equation}
    x_{\tau_{i-1}} = \sqrt{\bar{\alpha}_{\tau_{i-1}}}\,\hat{x}_0 + c\,\epsilon_\theta(x_{\tau_i},\tau_i) + \sigma_i z,
    \qquad
    \hat{x}_0 = \frac{x_{\tau_i} - \sqrt{1-\bar{\alpha}_{\tau_i}}\,\epsilon_\theta}{\sqrt{\bar{\alpha}_{\tau_i}}}
\end{equation}

其中的噪声\textbf{不是推导出来的，而是巧妙设计出来的}，刻意拆成两股，原因有三：\textbf{顺应代数结构}（$\hat{x}_0$ 代回后，展开项天然掉出一个与 $\epsilon_\theta$ 成比例的系数，写成 $c\,\epsilon_\theta$ 便于合并同类项）；\textbf{方差预算的"切蛋糕"逻辑}——总预算 $1-\bar{\alpha}_{\tau_{i-1}}$ 固定，手里恰有两股独立噪声来源（历史方向 $\epsilon_\theta$ 保持轮廓连续，新噪声 $z$ 弥补预测误差、增加多样性），独立噪声的方差具有可加性，正好切分；\textbf{保持高斯性}（独立高斯的线性组合仍是高斯，只需跟踪均值与方差，避免数学复杂化）。灵魂比喻：$c\,\epsilon_\theta$ 是手里的\textbf{指南针}（确定性引导），$\sigma_i z$ 是随机迈出的\textbf{小碎步}（随机性探索），两者相加就是"方向引导"与"随机探索"之间最简单的平衡。

\textbf{4. 系数确定与 $\eta$ 的调节}。由方差预算 $c^2+\sigma_i^2 = 1-\bar{\alpha}_{\tau_{i-1}}$ 解出 $c = \sqrt{1-\bar{\alpha}_{\tau_{i-1}}-\sigma_i^2}$：一旦选定 $\sigma_i$，$c$ 被唯一确定，而 $\sigma_i$ 本身是自由的——这正呼应"训练只约束边缘分布，转移规则不唯一"的观察。以 DDPM 广义后验标准差为基准引入 $\eta\in[0,1]$，令 $\sigma_i = \eta\,\sigma_i^{DDPM}$，代回得统一的 DDIM 反向更新公式：

\begin{equation}
    \boxed{x_{\tau_{i-1}} = \sqrt{\bar{\alpha}_{\tau_{i-1}}}\,\hat{x}_0 + \sqrt{1-\bar{\alpha}_{\tau_{i-1}}-\sigma_i^2}\;\epsilon_\theta(x_{\tau_i},\tau_i) + \sigma_i z}
\end{equation}

$\eta=1$ 时整式退化为 DDPM 的 ancestral sampling（子序列取完整时即标准 DDPM 公式）；$\eta=0$ 时 $\sigma_i=0$，无新噪声注入，过程完全确定——相同 $x_T$ 必得相同结果，"implicit" 由此得名。

\textbf{5. 少步采样为何倾向 $\eta=0$ 的确定性模式}。步数 $S$ 越小，子序列相邻步的噪声级差越大，$\hat{x}_0$ 的预测误差被放大得越多；$\eta=1$ 每步还要额外注入新噪声，相当于在已偏离的路径上再"乱抖一下"，误差逐级累积，样本明显模糊；$\eta=0$ 则完全消除了随机扰动这一误差来源。直观理解：步子跨得越大（步数越少），越要攥紧指南针（$\eta=0$），少瞎蹦跶；步数足够多（$S$ 接近 $T$）时偶尔掷骰子（$\eta=1$）才不至于跑偏，此时两者差别不大。需要注意，这一经典结论在本实验中\textbf{并非无条件成立}：余弦调度下 $\eta=0$ 采样曾全面失效，其定位与修复见 \ref{sec:eta0-failure} 节。

\subsection{扩展二：指定数字生成}

无条件模型生成哪个数字全凭运气；希望像给画家下订单一样指定数字 $y\in\{0,...,9\}$。标签先经可学习查找表变成向量 $\vec{emb}_y = \mathrm{EmbeddingTable}[y]$，与时间嵌入同维相加，模型变为 $\epsilon_\theta(x_t,t,y)$，而加噪、损失、采样公式全部不变。

只注入条件，模型的服从往往不够"坚决"。由贝叶斯公式，$\nabla_{x_t}\log p(x_t|y) = \nabla_{x_t}\log p(x_t) + \nabla_{x_t}\log p(y|x_t)$：条件分布的 score = 无条件 score + 分类器方向，给后者加权重 $w$ 即可调节条件强度。但不必真的训练分类器~\parencite{ho2022classifierfree}：$\nabla\log p(y|x_t) = \nabla\log p(x_t|y) - \nabla\log p(x_t)$，右边两项恰是同一网络在有/无条件下的两种输出：训练时以 10\% 概率把标签替换为空标签 $\varnothing$（embedding 表第 11 行），同一网络便同时学会两种模式。噪声预测与 score 只相差一个负系数~\parencite{song2021score}，代入得每步实际使用的引导预测：

\begin{equation}
    \boxed{\hat{\epsilon} = \epsilon_\theta(x_t,t,\varnothing) + w\left[\epsilon_\theta(x_t,t,y) - \epsilon_\theta(x_t,t,\varnothing)\right]}
\end{equation}

$w=1$ 退化为普通条件采样；$w>1$ 条件特征更鲜明（过大导致失真）；$w=0$ 退化为无条件采样。每步需条件/无条件两次前向，之后 DDPM/DDIM 公式照常使用。

\section{实验结果}

\subsection{实验设置}

\begin{table}[h]
    \centering
    \caption{模型与训练超参数}
    \label{tab:hyperparams}
    \begin{tabular}{l l}
        \toprule
        \textbf{项目} & \textbf{值} \\
        \midrule
        模型参数量 & 1.85M \\
        扩散步数 $T$ & 1000 \\
        噪声调度 & 余弦（本报告结果）；线性（可切换） \\
        优化器 & Adam，lr $= 10^{-3}$ \\
        batch size / epochs & 128 / 20 \\
        EMA 衰减 $\rho$ & 0.999 \\
        梯度裁剪 & clip\_grad\_norm $=1.0$ \\
        标签 dropout& 0.1 \\
        运行环境 & torch 2.14+CPU \\
        \bottomrule
    \end{tabular}
\end{table}

\subsection{损失曲线}

\begin{figure}[h]
    \centering
    \placeholderfig{figs/loss_curve.png}{0.7}
    \caption{训练损失曲线（每 epoch 平均 MSE）}
    \label{fig:loss}
\end{figure}

训练共 20 epoch（余弦调度，条件模型），平均 MSE 从首个 epoch 的 0.0855 快速下降到第 5 epoch 的约 0.040，随后进入缓慢收敛的平台期，最终稳定在 \textbf{0.0373}。曲线全程单调下降、无异常抖动，表明 lr $=10^{-3}$ 与梯度裁剪的配置合适；平台期的缓慢下降主要来自 EMA 之外参数的细微调整，属于小模型在小数据集上的正常形态。

\subsection{样本网格}

训练结束后从 $\mathcal{N}(0,I)$ 采 64 张纯噪声，走满 1000 步 ancestral sampling，拼成 8$\times$8 网格：

\begin{figure}[h]
    \centering
    \placeholderfig{figs/samples\_grid.png}{0.65}
    \caption{最终模型的 8$\times$8 样本网格（ancestral sampling，1000 步，EMA 权重）}
    \label{fig:grid}
\end{figure}

64 张样本中约 55$\sim$58 张（约九成）结构完整、可清晰辨认，数字类别分布均匀、笔画风格多样。典型失败形态为个别样本笔画断裂或与背景噪点粘连成团（图中可见 5$\sim$6 张），未出现整张纯噪声或模式坍塌（全部生成同一数字）的情况。对比精简配置（1 块/层、无 EMA）时期约 10/64 难辨认的结果，恢复容量并加入 EMA 后质量提升明显。

\subsection{不同采样步数对比（扩展一）}

实验协议：固定随机种子，使各步数下初始噪声 $x_T$ 完全相同，质量差异只来自步数与 $\eta$；步数取 $S \in \{1000, 200, 100, 50, 20\}$，分别用 $\eta=1$（DDPM 后验）与 $\eta=0$（确定性 DDIM）采样 64 张并记录耗时。

\begin{figure}[h]
    \centering
    \placeholderfig{figs/steps\_compare\_eta1.png}{0.9}
    \caption{$\eta=1$ 时不同采样步数的样本对比（每行一个步数，相同初始噪声）}
    \label{fig:steps-eta1}
\end{figure}

\begin{figure}[h]
    \centering
    \placeholderfig{figs/steps\_compare\_eta0.png}{0.9}
    \caption{$\eta=0$ 时不同采样步数的样本对比（每行一个步数，相同初始噪声）}
    \label{fig:steps-eta0}
\end{figure}

\begin{table}[h]
    \centering
    \caption{不同采样步数的耗时（64 张，CPU）}
    \label{tab:steps-time}
    \begin{tabular}{l l l}
        \toprule
        \textbf{步数 $S$} & \textbf{总耗时 (s)} & \textbf{平均每步 (ms)} \\
        \midrule
        1000 & 93.9 & 93.9 \\
        200 & 27.8 & 138.9 \\
        100 & 15.4 & 154.3 \\
        50 & 7.1 & 141.9 \\
        20 & 2.8 & 139.8 \\
        \bottomrule
    \end{tabular}
\end{table}

\textbf{耗时}：与步数近似线性（每步一次网络前向，约 100$\sim$140ms/步），从 1000 步的 93.9s 降到 20 步的 2.8s，加速约 34 倍；$\eta$ 不改变每步计算量，两条曲线耗时一致。

\textbf{质量}：$\eta=1$ 在 1000/200 步质量良好，$\leq$100 步出现明显噪声残留与轮廓模糊，20 步几乎只剩虚影——每步注入的新噪声在少步时逐级累积误差。$\eta=0$（经 \ref{sec:eta0-failure} 节所述的 $\hat{x}_0$ 截断修复后）表现截然不同：1000 到 20 步的样本几乎无差别，20 步仍清晰可辨，验证了 \ref{sec:ddim} 节"少步倾向确定性模式"的分析。也就是说，DDIM 在本实验中以 1/50 的步数（20 步、2.8s）达到了与 1000 步 ancestral sampling（93.9s）目测相当的质量。

复现命令见项目 \texttt{README.md} 的 Quick Start 一节。

\subsection{指定数字采样（扩展二）}

用条件模型分别以 $w=2.0$ 生成 0--9 各 8 张（每行一个数字），并对比不同 guidance 强度的效果：

\begin{figure}[h]
    \centering
    \placeholderfig{figs/digit\_all.png}{0.8}
    \caption{指定数字采样：0--9 各 8 张，每行一个数字（$w=2.0$，EMA 权重）}
    \label{fig:digit-all}
\end{figure}

\begin{figure}[h]
    \centering
    \placeholderfig{figs/digit\_w\_compare.png}{0.8}
    \caption{不同 guidance 强度 $w \in \{0, 1, 2, 4\}$ 的对比（同一数字、相同初始噪声）}
    \label{fig:digit-w}
\end{figure}

指定数字采样效果良好：图 \ref{fig:digit-all} 中 10 行样本全部命中指定数字，目测命中率约 100\%，且每行内部保留了笔画粗细、倾斜角度的多样性（非复制同一模板）。guidance 强度对比（图 \ref{fig:digit-w}，数字 8）显示：$w=0$ 退化为无条件采样（数字混杂）；$w=1$ 已稳定生成 8；$w=2$ 形态更饱满、一致性更高；$w=4$ 笔画进一步加粗、个别样本略显夸张，但未出现失真崩坏。$w\in[1,4]$ 均有效，$w=2.0$ 是条件强度与自然度之间的合适折中。

复现命令见项目 \texttt{README.md} 的 Quick Start 一节。

\section{遇到的问题与解决方法}

\subsection{训练耗时约为预期的 6 倍：参数量估算代替不了算力估算}

\textbf{现象}：初版模型（base 32、每层级 2 残差块、上采样带平滑卷积，3.16M 参数）按"3M 参数很小"预估 0.5$\sim$1h 训完，实测每 epoch 约 18 min，20 epoch 需约 6h。

\textbf{定位}：编写 \texttt{benchmark.py} 实测：单步训练约 230 GFLOPs（前向 0.60 GFLOPs/图 $\times$ 3 $\times$ batch 128），CPU 实际吞吐约 100 GFLOPS $\rightarrow$ 2.3 s/步，机器效率正常，瓶颈是纯卷积算力（前向+反向占 66\% 时间；数据加载仅 1.4\%）。MACs 分布显示：解码器三层占 43\%（跳跃连接拼接后 $2C \rightarrow C$ 的卷积最贵），两个上采样平滑卷积占 19\%，而 28$\times$28 分辨率处一次 $3\times3$ 卷积的成本是 7$\times$7 处的 16 倍。

\textbf{解决}：按实测收益排序应用四项优化——channels\_last 内存格式（$-25\%$ 时间）、base 32$\rightarrow$24（$-44\%$ MACs）、2 块 $\rightarrow$ 1 块每层（$-29\%$ MACs）、去掉上采样平滑卷积（$-19\%$ MACs），epochs 20$\rightarrow$15。单步 2.3s $\rightarrow$ 0.9s，15 epoch 约 40 min 训完。

\begin{lightnote}
教训：$\leq$5M 约束的是\textbf{参数量}而非\textbf{算力}。小参数模型若把通道堆在高分辨率层，FLOPs 依然巨大；估算训练时间必须分析每层的 $H \times W \times C^2$ 而不是参数总量。
\end{lightnote}

\subsection{精简模型样本质量不足：质量与预算的再平衡}

\textbf{现象}：为解决耗时问题精简后的 1.23M 参数模型（1 块/层），64 张样本中约 10 张难以辨认。

\textbf{分析与解决}：容量与训练量均不足。在 2h 预算的剩余空间内做三件事：恢复每层级 2 个残差块（1.23M $\rightarrow$ 1.85M，各分辨率容量翻倍）；新增 EMA 权重平均（每步开销 $<$1ms，对采样稳定性提升最直接）；epochs 15 $\rightarrow$ 20。不恢复 base 32 与平滑卷积（分别 $+76\%$、$+19\%$ 算力，性价比低）。最终配置 20 epoch 全程控制在预算内，样本质量满足实验要求。

\subsection{\texorpdfstring{$\eta=0$}{eta=0} 采样全面失效：余弦调度末端的首步发散}
\label{sec:eta0-failure}

\textbf{现象}：步数对比实验中，$\eta=0$（确定性 DDIM）在全部步数档位（含 1000 步）产出的样本都是无法辨认的噪声团，而 $\eta=1$ 完全正常——与"少步时 $\eta=0$ 更好"的经典预期恰好相反。

\textbf{定位}：对采样轨迹插桩，逐一步记录 $|x|$ 与 $|\hat{x}_0|$ 的量级：

\begin{table}[h]
    \centering
    \caption{50 步采样的轨迹量级（前 3 步与末步，$\max|x|$ / $\max|\hat{x}_0|$）}
    \label{tab:eta0-diag}
    \begin{tabular}{l l l l l}
        \toprule
        & \multicolumn{2}{c}{$\eta=0$} & \multicolumn{2}{c}{$\eta=1$} \\
        \cmidrule(lr){2-3} \cmidrule(lr){4-5}
        步 & $|x|$ & $|\hat{x}_0|$ & $|x|$ & $|\hat{x}_0|$ \\
        \midrule
        第 1 步（$t=999$） & 3.8 & \textbf{839.9} & 3.8 & \textbf{839.9} \\
        第 2 步（$t=979$） & 30.0 & 151.3 & 24.8 & 128.4 \\
        第 3 步（$t=958$） & 34.9 & 89.5 & 17.5 & 36.7 \\
        $\vdots$ & $\vdots$ & $\vdots$ & $\vdots$ & $\vdots$ \\
        最末步（$t=0$） & \textbf{26.4} & 26.1 & \textbf{1.04} & 1.02 \\
        \bottomrule
    \end{tabular}
\end{table}

数据范围本应是 $[-1,1]$，但 $\eta=0$ 的 $|x|$ 从第一步起持续发散、全程未恢复；$\eta=1$ 起点完全相同，却在两三步内被拉回正常量级。

\textbf{原因}：余弦调度在 $t \to T$ 时把 $\bar{\alpha}_t$ 压到接近 0（本模型 $\bar{\alpha}_T \approx 4\times10^{-33}$，而线性调度约 $4\times10^{-5}$）。构造 $\hat{x}_0 = (x_t - \sqrt{1-\bar{\alpha}_t}\,\epsilon_\theta)/\sqrt{\bar{\alpha}_t}$ 时要除以 $\sqrt{\bar{\alpha}_T} \approx 6\times10^{-17}$——网络预测与 $x_T$ 之间的任何微小偏差都被放大到数百倍（$|\hat{x}_0|$ 达 840）。$\eta=0$ 不注入任何新噪声，误差在确定性轨迹上步步相传、无法自恢复；而 $\eta=1$ 在高噪声区的后验方差 $\sigma^2 \approx 1$，注入的新鲜噪声占主导，相当于每步"重新洗牌"，一步之内就冲掉了放大的误差。换言之，\textbf{确定性方法没有随机性做掩护，数值细节会被放大为成败之别}；余弦调度把 $\bar{\alpha}$ 末端压得过低，本就不适合跳步/确定性采样（Nichol 与 Dhariwal 也指出线性调度更适合少步采样~\parencite{nichol2021improved}）。

\textbf{解决}：将 $\hat{x}_0$ 截断到数据范围 $[-1,1]$（DDIM 官方实现中的标准细节），一行改动。修复后 $\eta=0$ 全程 $|x| \leq 1$，且在 20$\sim$1000 步间质量几乎不变（见 \ref{fig:steps-eta0} 与 3.4 节分析）。

\begin{lightnote}
教训：反解公式中"除以接近 0 的量"是数值地雷。随机性采样器的噪声会掩盖这类问题，使其长期潜伏；一旦换用确定性采样器便立即暴露。实现论文公式时，应同步对照官方实现中的工程细节（截断、截位、eps 保护），它们往往不是可有可无的。
\end{lightnote}

\newpage

% ------------------------------------------- %
\printbibliography[heading=bibintoc, title={References}]
% ------------------------------------------- %

\end{document}
